% Vision-Language-Action (VLA) models aim to provide a single generalist controller for robots, but today's systems fall short on the criteria that matter for real-world deployment. Frontier models are closed, open-weight alternatives are tied to expensive hardware, reasoning-augmented policies pay prohibitive latency for their grounding, and fine-tuned success rates remain below the threshold for dependable use.
% We present \molmoactt, a fully open action reasoning model built for practical deployment, advancing over its predecessor along five axes. We introduce \molmoer, a VLM backbone specialized for spatial and embodied reasoning, trained on a 3.3M-sample spatial and embodied corpus with a specialize-then-rehearse recipe. We release three new datasets targeting platforms across the low-to-medium cost range: \molmoacttyamdata, 720 hours of teleoperated bimanual trajectories that constitute the largest open bimanual dataset to date; \molmoacttdroiddata, a quality-filtered Franka subset of DROID; and \molmoacttsodata, a quality-filtered SO-100/101 subset. We provide \molmoactfast, an open-weight, open-data action tokenizer trained on millions of trajectories across five embodiments. We redesign the architecture to graft a flow-matching continuous-action expert onto a discrete-token VLM via per-layer key-value (KV) conditioning. Finally, we propose \molmothink, an adaptive-depth reasoning variant that re-predicts depth tokens only for scene regions that change between timesteps, retaining geometric grounding at a fraction of prior latency.
% In the most extensive empirical study of any open VLA to date, spanning 7 simulation and real-world benchmarks, \molmoactt outperforms strong baselines including $\pi_{0.5}$, while \molmoer surpasses GPT-5 and Gemini Robotics ER-1.5 across 13 embodied-reasoning benchmarks. We release model weights, training code, and the complete training data, moving the open-source VLA stack meaningfully closer to one that researchers and practitioners can both build on and deploy.


Vision-Language-Action (VLA) models aim to provide a single generalist controller for robots, but today's systems fall short for real-world deployment. Frontier models are closed, open-weight alternatives are tied to expensive hardware, reasoning-augmented policies pay prohibitive latency for their grounding, and fine-tuned success rates remain below the threshold for dependable use. We present \molmoactt, a fully open action reasoning model built for practical deployment, advancing its predecessor, \molmoact along five axes. We introduce \molmoer, a VLM backbone specialized for spatial and embodied reasoning, trained on a 3.3M-sample corpus with a specialize-then-rehearse recipe. We release three new datasets spanning low-to-medium cost platforms: \molmoacttyamdata, 720 hours of teleoperated bimanual trajectories that constitute the largest open bimanual dataset to date; \molmoacttdroiddata, a quality-filtered Franka subset of DROID; and \molmoacttsodata, a quality-filtered SO-100/101 subset. We provide \molmoactfast, an open-weight, open-data action tokenizer trained on millions of trajectories across five embodiments. We redesign the architecture to graft a flow-matching continuous-action expert onto a discrete-token VLM via per-layer key-value (KV) conditioning. Finally, we propose \molmothink, an adaptive-depth reasoning variant that re-predicts depth tokens only for scene regions that change between timesteps, retaining geometric grounding at a fraction of prior latency. In the most extensive empirical study of any open VLA to date, spanning 7 simulation and real-world benchmarks, \molmoactt outperforms strong baselines including $\pi_{0.5}$, while \molmoer surpasses GPT-5 and Gemini Robotics ER-1.5 across 13 embodied-reasoning benchmarks. We release model weights, training code, and complete training data.
